
\section{Introduction} \label{sec:intro}


Graph streaming algorithms process their inputs by making one or few passes
over the edges of an input graph using limited memory.
Algorithms that use space proportional to $n$, the number of vertices, are
called \emph{semi-streaming} algorithms.
Since their introduction by~\cite{FeigenbaumKMSZ05}, graph streaming algorithms
have become one of the main theoretical research areas on  processing massive
graphs.
We refer the interested
reader to~\cite{McGregor14} for an introductory survey of earlier results on
this topic. 

In this work, we prove a polynomial-pass lower bound for any graph streaming algorithm
that computes $k$-cores or degeneracy of a given graph.  Our result is of interest from the point of view of proving strong lower bounds in the graph streaming model in addition to their direct implications 
for these two specific problems. 

\subsection{Polynomial Pass Lower Bounds in Graph Streams}\label{sec:poly-pass}

Even though the study of multi-pass graph streaming algorithms started hand in hand with single-pass algorithms in~\cite{FeigenbaumKMSZ05}, our understanding of 
powers and limitations of multi-pass algorithms, even for most basic problems, lags considerably behind. On one hand, for a problem like minimum cut, 
we have algorithms that in just $\Ot{n}$ space and two passes can solve the problem \emph{exactly}~\cite{AssadiD21}\footnote{See~\cite{RubinsteinSW18} for an implicit algorithm with the same bounds and~\cite{MukhopadhyayN20} for the extension to 
weighted cuts in $O(\log{n})$ passes.} (see~\cite[Table 1]{AssadiCK19} for a list of several such results). On the other hand, for some other basic problems such as undirected shortest path, directed reachability, and bipartite matching, the best known semi-streaming algorithms require $O({n^{1/2}})$~\cite{ChangFHT20}, 
$n^{1/2+o(1)}$~\cite{LiuJS19,AssadiJJST22}, and $n^{3/4+o(1)}$~\cite{AssadiJJST22} passes, respectively; yet, despite significant efforts, the best lower bound for any of these problems
is still (even slightly below) $\Omega(\log{n})$ passes~\cite{GuruswamiO13,ChakrabartiGMV20,AssadiR20,ChenKPSSY21}. 

A key reason behind our  weaker understanding of multi-pass streaming algorithms can be attributed to the lack of techniques for 
proving  \emph{super-logarithmic} pass lower bounds for semi-streaming algorithms. At this point, such lower bounds
are only known for a handful of problems: clique and independent set~\cite{HalldorssonSSW12}, dominating set~\cite{Assadi17sc}, Hamiltonian path~\cite{BachrachCDELP19}, maximum cut~\cite{BachrachCDELP19,KolPSY23}, vertex cover and coloring~\cite{AbboudCKP21}, exact Boolean CSPs~\cite{KolPSY23}, triangle detection~\cite{PapadimitriouS84,Bar-YossefKS02}, and diameter computation~\cite{FrischknechtHW12}. Although, for all these problems, we can 
actually prove  close-to-$n$ pass lower bounds. Let us examine this dichotomy. 

A quick glance at the list of problems above may suggest an intuitive difference between these problems and the ones like reachability or shortest path: the above list consists of problems 
that are computationally hard in a classical sense\footnote{These are  standard NP-hard problems or admit some fine-grained hardness (for the latter two)~\cite{RodittyW13,KopelowitzPP16}.}, 
suggesting that we are dealing with a ``harder'' class of problems in their case. While this intuition should not be taken as a formal evidence---as classical computational hardness does \emph{not} imply 
streaming lower bounds (which are unconditional and information-theoretic)---\cite{AssadiCK19} showed that one can also formally explain this dichotomy.

In particular, \cite{AssadiCK19} observed that these strong streaming lower bounds happen only when the communication complexity of the problem at hand
 is $\Omega(n^2)$. Such a high lower-bound on the communication complexity \textit{immediately} gives an $\Omgt{n}$-pass lower bound for semi-streaming algorithms via standard reductions. Whereas, for almost all problems of interest in the semi-streaming model, including shortest path, reachability, and bipartite matching, 
we already know an $\Ot{n}$ communication upper bound\footnote{This perhaps can be seen as this: a problem whose (unbounded round) communication complexity is already high has almost no place in the streaming
model, which is a much weaker model algorithmically.} (the protocol for bipartite matching was only discovered in~\cite{blikstadBEMN22} after the work of~\cite{AssadiCK19}, but $\Ot{n^{3/2}}$ communication protocols were known already~\cite{IvanyosKLSW12,DobzinskiNO19}). We refer the reader to~\cite[Section 1.1]{AssadiCK19} for more context regarding these observations and prior techniques for $o(\log{n})$ pass lower bounds. 

\subsubsection*{Toward Stronger Streaming Lower Bounds} 
A natural question in light of these observations, already posed in~\cite{AssadiCK19}, is the following: 

\begin{quote}
\textbf{Motivating question.}
\textit{Is there any graph problem which is ``not too hard'', in that it can be solved efficiently with communication (nearly) linear in the number $n$ of vertices, and for which, nonetheless, any semi-streaming algorithm needs a polynomial $n^{\Omega(1)}$ number of passes?}
\end{quote}
To address this question,~\cite{AssadiCK19} introduced a new (four-player) communication problem called \textbf{Hidden Pointer Chasing (HPC)}, 
which acts as a cross between \emph{Set-Intersection} and \emph{Pointer Chasing} problems, which are the main problems for, respectively, proving $\Omega(n^2)$ communication lower bounds on graphs, and $o(\log{n})$-pass lower
bounds for semi-streaming algorithms. 

Roughly speaking, the HPC problem is defined as follows. There are four players paired into two groups. Each pair
of players inside a group shares $m$ instances of the Set-Intersection problem on $m$ elements (each of the two players holds a subset of $[m]$ and they need to identify the unique intersecting element). The
intersecting element in each instance of each group ``points'' to an instance in the other group. The
goal is to start from a fixed instance, follow these pointers for a fixed number of steps, and then return the last element reached. See~\Cref{prob:hpc} for the formal description. 

This problem admits an efficient communication protocol with no limit on its number of rounds, but~\cite{AssadiCK19} showed that any $r$-round protocol
that aims to find the $(r+1)$-th pointer in HPC requires $\Omega(m^2/r^2)$ communication. This places HPC 
squarely in the middle of previous techniques and quite suitable for performing reductions to prove streaming lower bounds
even for not-too-hard graph problems. Using this,~\cite{AssadiCK19} proved the first set of polynomial-pass graph streaming lower bounds
in this class of problems: computing \emph{lexicographically-first maximal independent set (LFMIS)} and \emph{$s$-$t$ minimum cut} on graphs with exponential edge-capacities both require $\Omgt{n^{1/5}}$ passes to be solved by semi-streaming algorithms. 

Despite the significant advances on multi-pass streaming lower bounds in the last couple of years~\cite{ChakrabartiGMV20,AssadiKSY20,AssadiR20,ChenKPSSY21,AssadiN21,ChenKPSSY21b,Assadi22,AssadiKZ22,KolPSY23,ChenKPSSY23}, 
there is still no other known (not-too-hard) problem that admits a polynomial-pass lower bounds beside those of~\cite{AssadiCK19}. In addition, it is worth mentioning that, strictly speaking, neither LFMIS nor the version of $s$-$t$ minimum cut in~\cite{AssadiCK19} 
completely fit the premise of our original question: LFMIS is not purely a graph problem as it is not invariant under labeling of the vertices, and $s$-$t$ minimum cut  studied in~\cite{AssadiCK19} involves (i) making 
the non-standard assumption of exponential capacities, and (ii) even for unit-capacity graphs, is not known to admit an $\Ot{n}$ communication protocol (see~\cite{blikstadBEMN22}). 

\medskip
\noindent
We prove polynomial-pass lower bounds for two natural graph problems, $k$-cores and degeneracy, by reduction from a harder variant of the HPC problem which we call MultiHPC. We also present novel $\tO(n)$ communication protocols for these two problems. These two results together give us the first natural instances of a positive answer to our motivating question. 

These results further demonstrate the power of reductions from the HPC problem, as a technique for proving strong lower bounds in the graph streaming model, which are beyond the reach of other techniques. With this in mind, we improve the lower bound of~\cite{AssadiCK19} for the HPC problem to an optimal bound of $\Omega(m^2/r)$ communication, i.e., we improve the known bound by a factor of $r$. This contribution alone results in a polynomial improvement in the number of passes, for all lower bounds that follow via reductions from HPC (for instance, it immediately improves the bounds of~\cite{AssadiCK19} for LFMIS and exponential-capacity $s$-$t$ minimum cut to $\Omgt{n^{1/4}}$ passes).

But, as it turns out, all the known lower-bounds that follow by reduction from HPC also follow by reduction from MultiHPC. For this variant, we can prove an $\Omega(m^2)$ lower-bound for $r$ rounds (since the input size for MultiHPC is $r\cdot m^2$), and this translates to an improved semi-streaming lower-bound of $\Omgt{n^{1/3}}$ passes for all of the above problems.




\subsection{\texorpdfstring{$k$-Cores}{k-Cores} and Degeneracy in Graph Streams}\label{sec:degen-intro}

For any undirected graph $G=(V,E)$ and integer $k \geq 1$, a \emph{$k$-core} in $G$ is a maximal set $S$ of vertices such that the induced subgraph of $G$ on $S$, denoted by $G[S]$, has a minimum degree of at least $k$. In other words, 
any vertex in $S$ has at least $k$ other neighbors in $S$. 

$k$-Cores provide a natural notion of well-connectedness in massive graphs, and as such, computing $k$-cores (and more generally $k$-core decompositions; see, e.g.,~\cite{LiuSYDS22}) 
has been widely studied in databases~\cite{BonchiGKV14, ChuZ00ZXZ20, LiZZQZL19}, social networks~\cite{DhulipalaBS17, DhulipalaBS18, KhaouidBST15}, machine learning~\cite{Alvarez-HamelinDBV05, EsfandiariLM18, GhaffariLM19}, 
among others~\cite{LiWSX18,GalimbertiBGL20, SariyuceGJWC13}. 

As a result, in recent years, there has been a rapidly growing body of work on
computing $k$-cores on massive graphs in parallel and streaming
models of computation~\cite{SariyuceGJWC13,DhulipalaBS17, DhulipalaBS18,EsfandiariLM18, GhaffariLM19,LiuSYDS22}. In particular,~\cite{EsfandiariLM18} presented a single-pass algorithm that for any $\eps > 0$, 
computes a $(1-\eps)$-approximation of every $k$-core in $G$ (i.e., obtains a $(1-\eps)$-approximate $k$-core decomposition) in $\Ot{n/\eps^2}$ space (see also~\cite{SariyuceGJWC13} for an earlier streaming algorithm and~\cite{GhaffariLM19}
for a closely related parallel algorithm). 

The \emph{degeneracy} of a graph $G=(V,E)$, denoted by $\kappa(G)$, is the largest integer $k \geq 0$ such that $G$ contains a non-empty $k$-core. 
The simple greedy algorithm that at every step peels off the smallest degree vertex results in the so-called \emph{degeneracy ordering} of $G$ and $\kappa{(G)}$ 
is equal to the largest degree of a vertex removed in this peeling process~\cite{MatulaB83}. Degeneracy is a standard measure of uniform sparsity and is closely
related to other such notions like arboricity (which is within a factor $2$ of degeneracy). Moreover, computing degeneracy is a subroutine for approximating various
other problems such as arboricity~\cite{AlonYZ97}, densest subgraph~\cite{Charikar00}, $(\kappa+1)$ coloring~\cite{erdHos1966chromatic}. 

The degeneracy problem, and closely related uniform-sparsity measures such as densest subgraph, have also been studied extensively in the graph 
streaming literature~\cite{BahmaniKV12,Farach-ColtonT14,BhattacharyaHNT15,McGregorTVV15,Farach-ColtonT16,BeraCG19,AlonA20}. In particular,~\cite{Farach-ColtonT14} provided an $O(\log{n})$-pass semi-streaming algorithm that outputs a constant factor approximation to degeneracy, and~\cite{Farach-ColtonT16} subsequently improved this to a single-pass $(1-\eps)$-approximation in $\Ot{n/\eps^2}$ space (see also~\cite{McGregorTVV15} for densest subgraph and~\cite{BeraCG19,AlonA20} for degeneracy coloring). 
 
\medskip

In terms of lower bounds, \cite{BeraCG19} prove that any \emph{single-pass} streaming algorithm that computes the exact value of degeneracy or approximates it within an additive factor of $\lambda$ requires $\Omega(n^2)$ space or $\Omega(n^2/\lambda^2)$ space respectively. Our \emph{polynomial-pass} lower bounds for $k$-cores and degeneracy, now in a very strong sense, rule out the possibility of extending any prior semi-streaming algorithms computing near-optimal solutions to these problems, to compute \emph{exactly} optimal solutions. 

\subsection{Our Results}\label{sec:res}

We give an informal presentation of our results here and postpone the formal descriptions to subsequent sections. The first main result is 
our polynomial-pass lower bound for $k$-core computation and degeneracy. 
\\

\begin{result}[Formalized in~\Cref{thm:degen-lb}]\label{res:stream-lb}
	For any integer $p \geq 1$, any $p$-pass streaming algorithm for computing the degeneracy of an input $n$-vertex graph requires $\Omgt{n^2/p^3}$ space. In particular, any semi-streaming algorithm for the problem requires $\Omgt{n^{1/3}}$ passes. 

\noindent	
	Moreover, the same lower bounds also apply 
	to the algorithms that given any integer $k \geq 1$, can check whether or not the input graph contains a non-empty $k$-core. 
\end{result}

\Cref{res:stream-lb} provides a strongly negative answer to the question of obtaining semi-streaming algorithms for exact computation of degeneracy and $k$-cores in a small number of passes. 
We obtain~\Cref{res:stream-lb} via a detailed and technical reduction, presented in~\Cref{sec:reduction}, from a variant of the Hidden Pointer Chasing (HPC) problem of~\cite{AssadiCK19}, which we call Boolean Multilayer Hidden Pointer Chasing (BMHPC). 

In a standard HPC problem, we are given $m$ instances of $m$-bit Set-Intersection ($O(m^2)$ bits in total) and interpret the intersection point of each instance as pointing to a different instance among these $m$. We then wish to know the position we end up in after following $r+1$ pointers. In a Boolean variant, we only care to know the \textit{parity} of the position we end up in. In a Multilayer variant, we are given $r$ different layers, each layer with its own $m$ instances ($r m^2$ bits in total), where we think of the intersection points at each layer as pointing to some instance in the next layer, and wish to know where we end up in the last layer by following these pointers.

In addition to the reduction from BMHPC, in~\Cref{sec:communication-upper}, we present a novel and non-trivial communication protocol that finds the degeneracy ordering (and thus degeneracy itself) and non-empty $k$-cores for any given $k$, using only $\Ot{n}$ communication. This communication upper bound thus places the $k$-core and degeneracy problems as perfect illustrations of a positive answer to our and~\cite{AssadiCK19}'s motivating question outlined earlier: namely, problems 
that prior techniques could not have proven any lower bound beyond $\log{n}$ passes. 
\Cref{res:stream-lb} thus constitutes the first set of natural graph problems with polynomial pass lower bounds for semi-streaming algorithms. 

\medskip

Our second main contribution is providing optimal lower bounds for the HPC problem and all its variants (Boolean, Multilayer, and Boolean Multilayer). A communication lower-bound of $\Omega(\frac{m^2}{r^2})$ was previously known for $(r-1)$-round protocols computing the $r$-th pointer in a (single layer) HPC problem. We prove the following.

\begin{result}[Formalized in~\Cref{thm:lb-hpc}]\label{res:hpc-lb}
	For any integer $1 \le r = O(\sqrt m)$, any $(r-1)$-round protocol for computing the $r$-th pointer in the HPC problem on a universe of size $m$ requires $\tilde \Omega({m^2}/r)$ communication. 

\medskip
\noindent
 For any integer $r \geq 1$, any $(r-1)$-round protocol for computing the $r$-th pointer in the Multilayer HPC problem on a universe of size $m$ requires $\tilde \Omega(m^2)$ communication. 
 
	\medskip
	\noindent
	Moreover, the same lower bounds hold for the Boolean versions of the above. 
\end{result}

\Cref{res:hpc-lb}, by strengthening the lower bound of~\cite{AssadiCK19}, allows us to prove polynomially stronger bounds on the number of passes of semi-streaming algorithms via reductions from HPC. As it turns out, every known reduction from HPC \cite{AssadiCK19} can be easily converted to a reduction from MHPC. Our results thus imply improved pass lower bounds (from $\Omgt{n^{1/5}}$ to $\Omgt{n^{1/3}}$)  for semi-streaming algorithms solving these problems. We capture this in the next corollary.

\begin{corollary}\label{cor:lfmis-mincut}
For any integer $p \geq 1$, any $p$-pass streaming algorithm for the following problems on $n$-vertex graphs requires $\Omgt{n^2/p^3}$ space. In particular, any semi-streaming algorithm for these problems require $\Omgt{n^{1/3}}$ passes.
\begin{itemize}
    \item Computing the minimum s-t cut value in a weighted graph (with exponential edge capacities)
    \item Computing the lexicographically-first maximal independent set (LFMIS) of an undirected graph
\end{itemize}
\end{corollary}

We obtain~\Cref{res:hpc-lb} by following the elegant analysis of pointer chasing problems due to~\cite{Yehudayoff20} via the \emph{triangular discrimination distance} between distributions, as opposed to more standard measures 
such as KL-divergence and total variation distance typically used in this context. This in turn requires extending the notion of ``almost solving'' for the Set-Intersection problem introduced by~\cite{AssadiCK19} (and further refined in~\cite{AssadiR20}),
to the triangular discrimination distance: roughly speaking, this corresponds to proving a lower bound for communication protocols that, instead of finding the intersecting element, change its distribution slightly 
from uniform distribution. The analysis in~\cite{AssadiCK19} measured this change by total variation distance, but now we need to do so by triangular discrimination distance instead. Finally, we prove a nearly-optimal
lower bound on the communication-distance tradeoff for almost solving Set-Intersection in terms of the triangular discrimination distance. 



%\clearpage
